% !TeX root = . . /main. tex

\chapter{概率模型证明中学恒等式的方法与实例}
概率模型是用概率来描述和分析随机现象的数学模型. 通过定义随机事件、样本空间、概率分布以及相关的概率规则, 来刻画不确定事件发生的可能性. 

% 在概率模型中, 随机变量用于表示随机试验的结果, 通常分为离散型和连续型两种. 样本空间是随机试验中所有可能结果的集合. 概率分布规定了随机变量取不同值的概率规律. 
% 在概率模型中，随机变量用于表示随机试验的结果. 随机变量通常分为离散型和连续型两种. 样本空间是随机试验中所有可能结果的集合. 概率分布规定了随机变量取不同值的概率规律. 


用概率模型解决高中数学问题时, 常见概率模型主要有以下几种:

等可能事件概率模型: 等可能事件的特性是随机试验中每个基本事件发生的可能性相等. 等可能事件概率模型就使用了等可能事件的这一特点. 在这种概率模型下, 事件 \(A\) 发生的概率 \(P\left( A\right)\) 等于事件 \(A\) 包含的事件数除以样本空间 \(\Omega\) 包含的事件总数. 例如, 投掷一枚质地均匀的骰子,出现每个点数的概率都是 \(\frac{1}{6}\) ,这就是一个等可能事件概率模型. 

独立重复试验模型: 本模型指的是在相同条件下重复进行一系列相互独立的试验. 每次试验只有两种可能结果(成功或失败). 而且每次试验成功的概率 \(p\) 保持不变. 例如, 多次抛掷一枚质地均匀的硬币, 每次抛硬币正面朝上 (反面朝上) 的概率都是 \(\frac{1}{2}\) . 每次试验的结果相互独立,这就是独立重复试验模型. 

相互独立事件概率模型: 通过构造多个相互独立的事件来解决问题. 若事件 \(A\) 、 \(B\) 相互独立,则 \(P\left( {AB}\right)  = P\left( A\right) P\left( B\right)\) . 例如,同时抛掷两枚质地均匀的骰子,第一枚骰子的结果与第二枚骰子的结果相互独立, 可利用这种独立性来计算相关事件的概率. 


\section{概率模型适用原理与证明思路}
% \subsection{原理与适用类型}
% 构建等可能事件概率模型证明恒等式的原理基于等可能事件的概率定义. 在等可能事件中, 每个基本事件发生的概率相等, 通过设计合理的随机试验, 将恒等式中的数学元素与等可能事件的概率计算相关联. 对于一些涉及组合数、排列数以及具有一定对称性的恒等式, 常常可以运用这种方法进行证明. 在组合恒等式中, 组合数可以表示从若干元素中选取特定数量元素的组合方式数,而通过构造摸球、分配物品等等可能事件, 可以将组合数转化为事件发生的概率, 利用概率的性质和计算规则来证明恒等式. 
构建等可能事件概率模型证明恒等式的原理基于等可能事件的概率定义. 在等可能事件中，每个基本事件发生的概率相等，举个例子，袋中现有$2n$个球（除颜色外别的均相同），其中$n$个白球，$n$个黑球，这时我们随机摸取一个，摸到白球和摸到黑球的概率是等可能的. 摸球模型也是最常用的等可能事件概率模型. 这样，我们可以通过设计合理的随机试验，将恒等式中的数学元素与等可能事件的概率计算关联起来. 对于一些含有组合数、排列数或者说是有一定对称性的恒等式，常常可以运用这种方法进行证明. 在组合恒等式中，组合数可以表示从若干元素中选取特定数量元素的组合方式数，而通过构造摸球、分配物品等等可能事件，可以将组合数转化为事件发生的概率，利用概率的性质和计算规则来证明恒等式. 



独立重复试验概率模型具有独特的特点, 它是在相同条件下, 重复进行 \(n\) 次试验, 每次试验的结果只有两种, 即事件 $A$ 发生或不发生, 且每次试验中事件 \(A\) 发生的概率 \(p\) 保持不变, 各次试验相互独立. 这种模型的概率计算基于二项分布, 在 \(n\) 次独立重复试验中, 事件 \(A\) 恰好发生 \(k\) 次的概率为 \(P\left( {X = k}\right)  = {C}_{n}^{k}{p}^{k}{\left( 1 - p\right) }^{n - k}\) . 

利用独立重复试验概率模型证明恒等式时, 关键在于将恒等式与试验结果建立紧密联系. 需要深入分析恒等式的结构特征, 确定其中的数学量与独立重复试验中的参数相对应. 若恒等式中出现组合数 \({C}_{n}^{k}\) 以及与 \(p\) 、$(1 - p)$相关的幂次形式, 就可以尝试将其与独立重复试验中事件 \(A\) 发生 \(k\) 次的概率表达式相联系. 在证明
\({C}_{n}^{0}{p}^{0}{\left( 1 - p\right) }^{n} + {C}_{n}^{1}{p}^{1}{\left( 1 - p\right) }^{n - 1} + \ldots  + {C}_{n}^{n}{p}^{n}{\left( 1 - p\right) }^{0} = 1\) 时, 可将其看作是在 \(n\) 次独立重复试验中, 事件 \(A\) 发生 0 次、1 次、 \(\cdots n\) 次的概率之和, 由于这些情况涵盖了所有 \(P\left( D\right)  = {\left\lbrack  1 - P\left( M\right) \right\rbrack  }^{2}\) 可能的试验结果, 所以其概率之和为 1, 从而利用独立重复试验概率模型完成了恒等式的证明. 

\section{实例分析}
% \begin{example}\label{example:例4}
%   证明恒等式 \({\left( a + b\right) }^{2} = {a}^{2} + {2ab} + {b}^{2}\) . 
% \end{example}
% \begin{proof}
%   构造如下概率模型,假设有一个袋子,里面装有 \(a\) 个红球和 \(b\) 个蓝球. 从其中有放回地抽取两次球. 

% 因此,我们可以设事件 \(A\) 为第一次抽到红球且第二次抽到红球,就有

% \[
% P\left( A\right)  = \frac{a}{a + b} \times  \frac{a}{a + b} = \frac{{a}^{2}}{{\left( a + b\right) }^{2}};
% \]

% 设事件 \(B\) 为第一次抽到红球且第二次抽到蓝球,就有

% \[
% P\left( B\right)  = \frac{a}{a + b} \times  \frac{b}{a + b} = \frac{ab}{{\left( a + b\right) }^{2}};
% \]

% 设事件 \(C\) 为第一次抽到蓝球且第二次抽到红球,就有

% \[
% P\left( C\right)  = \frac{b}{a + b} \times  \frac{a}{a + b} = \frac{ab}{{\left( a + b\right) }^{2}};
% \]

% 设事件 \(D\) 为第一次抽到蓝球且第二次抽到蓝球,就有

% \[
% P\left( D\right)  = \frac{b}{a + b} \times  \frac{b}{a + b} = \frac{{b}^{2}}{{\left( a + b\right) }^{2}}. 
% \]

% 四个事件构成了从袋子里两次抽球的所有情况, 因此, 四个事件的概率之和即为 1,所以有 \(P\left( A\right)  + P\left( B\right)  + P\left( C\right)  + P\left( D\right)  = 1\) ,即是 \(\frac{{a}^{2}}{{\left( a + b\right) }^{2}} + \frac{ab}{{\left( a + b\right) }^{2}} + \frac{ab}{{\left( a + b\right) }^{2}} + \frac{{b}^{2}}{{\left( a + b\right) }^{2}} = 1\) , 等式两边同乘 \({\left( a + b\right) }^{2}\) ,可得 \({\left( a + b\right) }^{2} = {a}^{2} + {2ab} + {b}^{2}\) . 
% \end{proof}
\begin{example}
  证明恒等式 \[(a + b)^n = \sum_{i=0}^n C_n^i a^i b^{n-i},\]\(a, b, n\) 均为自然数。
\end{example}

\begin{proof}
  现袋中有 \(a\) 个黑球，\(b\) 个白球。从其中有放回地随机摸取两次球。我们设事件 \(A_i\) 为从袋中摸出 \(i\) 个黑球，\(i = 1,2,3\ldots n\)。于是我们有
\[
P(A_i) = C_{n}^{i} \left( \frac{a}{a + b} \right)^i \left( \frac{b}{a + b} \right)^{n - i}. 
\]
摸出黑球的数量不同时，即 \(i\) 不同时显然有 \(A_i \cap A_j = \varnothing\)（\(i \neq j\)）。因此，
\[
\sum_{i = 0}^{n} P(A_i) = 1,
\]
即有
\begin{align*}
P(A_0) + P(A_1) + P(A_2) + \ldots + P(A_n) &= C_{n}^{0} \left( \frac{a}{a + b} \right)^0 \left( \frac{b}{a + b} \right)^n \\
&\quad + C_{n}^{1} \left( \frac{a}{a + b} \right) \left( \frac{b}{a + b} \right)^{n - 1} \\
&\quad + \ldots + C_{n}^{n} \left( \frac{a}{a + b} \right)^n \left( \frac{b}{a + b} \right)^0 \\
&= \sum_{i = 0}^{n} C_{n}^{i} \left( \frac{a}{a + b} \right)^i \left( \frac{b}{a + b} \right)^{n - i} \\
&= \frac{1}{(a + b)^n} \sum_{i = 0}^{n} C_n^i a^i b^{n - i} \\
&= 1. 
\end{align*}
因此，\(P(A_i) = C_{n}^{i} \left( \frac{a}{a + b} \right)^i \left( \frac{b}{a + b} \right)^{n - i}\) 成立. 
\end{proof}


\begin{example}\label{example:例3. 2}
  证明组合加法公式 \[C_{n+1}^k = C_n^k + C_n^{k-1}. \]
\end{example}

\begin{proof}
现设袋中有 \(n\) 个白球，1 个黑球。现从中随机摸取 \(k\) 个球（\(1 \leq k \leq n + 1\)）。

记事件 \(A\) 为“摸取的球中有黑球”，事件 \(B\) 为“摸球的球中没有黑球”。显然有 \(P(A) + P(B) = 1\)。

由 \(P(A) = \frac{C_{1}^{1} C_{n}^{k - 1}}{C_{n + 1}^{k}}\)，\(P(B) = \frac{C_{n}^{k}}{C_{n + 1}^{k}}\) 可得
\[
\frac{C_{n}^{k - 1}}{C_{n + 1}^{k}} + \frac{C_{n}^{k}}{C_{n + 1}^{k}} = 1. 
\]
等式两边同乘 \(C_{n + 1}^{k}\) 即有
\[
C_{n + 1}^{k} = C_{n}^{k} + C_{n}^{k - 1}. 
\]
恒等式成立. 
\end{proof}

\begin{example}
  证明恒等式 \[\sum_{k=0}^n \binom{n}{k}^2 = \binom{2n}{n}. \]
\end{example}

\begin{proof}
设袋中有 \(n\) 个白球和 \(n\) 个黑球。记事件 \(A\) 为从袋中随机摸取 \(n\) 个球，其中至少 1 个是白球。

取到 \(k\) 个白球的概率
\[
p_k = \frac{C_{n}^{k} C_{n}^{n - k}}{C_{2n}^{n}} = \frac{\left( C_{n}^{k} \right)^2}{C_{2n}^{n}}. 
\]
因此就有
\[
P(A) = \frac{\left( C_{n}^{1} \right)^2}{C_{2n}^{n}} + \frac{\left( C_{n}^{2} \right)^2}{C_{2n}^{n}} + \ldots + \frac{\left( C_{n}^{n} \right)^2}{C_{2n}^{n}}. 
\]
从另外一个角度来讲，计算摸不到白球的概率也可以得到 \(P(A)\)
\[
P(A) = 1 - \frac{C_{n}^{n} C_{n}^{0}}{C_{2n}^{n}}. 
\]
于是有
\[
\frac{\left( C_{n}^{1} \right)^2}{C_{2n}^{n}} + \frac{\left( C_{n}^{2} \right)^2}{C_{2n}^{n}} + \ldots + \frac{\left( C_{n}^{n} \right)^2}{C_{2n}^{n}} = 1 - \frac{\left( C_{n}^{0} \right)^2}{C_{2n}^{n}}. 
\]
移项后有
\[
\frac{\left( C_{n}^{0} \right)^2}{C_{2n}^{n}} + \frac{\left( C_{n}^{1} \right)^2}{C_{2n}^{n}} + \frac{\left( C_{n}^{2} \right)^2}{C_{2n}^{n}} + \ldots + \frac{\left( C_{n}^{n} \right)^2}{C_{2n}^{n}} = 1. 
\]
两边同乘 \(C_{2n}^{n}\) 后即可得
\[
\sum_{k = 0}^{n} \left( C_{n}^{k} \right)^2 = C_{2n}^{n}. \qedhere
\]
\end{proof}

\begin{example}
  证明 \[C_n^0 C_m^k + C_n^1 C_m^{k-1} + \cdots + C_n^k C_m^0 = C_{n+m}^k. \]
\end{example}

\begin{proof}
设袋中装有 \(n\) 个白球，\(m\) 个黑球，从中随机摸取 \(k\) 个（\(k \leq m\) 且 \(k \leq n\)）。

记事件 \(A\) 为从中随机摸取得 \(k\) 个球中恰好有 \(i\) 个白球。事件 \(A\) 的概率为 \(P(A_i) = \frac{C_{n}^{i} C_{m}^{k - i}}{C_{n + m}^{k}}\)，并且又因为各 \(A_i\) 不相容，有 \(P(A_0) + P(A_1) + \ldots + P(A_k) = 1\)。

即有
\[
\frac{C_{n}^{0} C_{m}^{k}}{C_{n + m}^{k}} + \frac{C_{n}^{1} C_{m}^{k - 1}}{C_{n + m}^{k}} + \ldots + \frac{C_{n}^{k} C_{m}^{0}}{C_{n + m}^{k}} = 1,
\]
两边同乘 $C_{n + m}^{k}$ 即得证。
\end{proof}

\begin{example}
  证明 \[C_n^0 + \frac{1}{2} C_{n+1}^1 + \frac{1}{4} C_{n+2}^2 + \cdots + \frac{1}{2^n} C_{2n}^n = 2^n. \]
\end{example}

\begin{proof}
现在有两个相同的盒子，编号分别为 1 和 2，其中各放有 \(n\) 个乒乓球，每次打球时任在一盒中取出一个乒乓球，求发现一盒用光时另一盒有 \(r\) 个的概率。现在记取盒子 1 内乒乓球为事件 \(A\)，取盒子 2 内乒乓球为事件 \(B\)，就有 \(p(A) = p(B) = \frac{1}{2}\)。已知一盒用光时另一盒有 \(r\) 个，那么总共抽取了 \(2n - r\) 次。

现在假设抽完的盒子是盒子 1，可知 \(P_{2n - r}(n) = C_{2n - r}^{n} \left( \frac{1}{2} \right)^{2n - r}\)。
若 \(r\) 从 0 取到 \(n\) 时，由各个事件不同时发生，知
\[
P_{2n} + P_{2n - 1} + P_{2n - 2} + \ldots + P_{n} = 1,
\]
即
\[
C_{2n}^{n} \left( \frac{1}{2} \right)^{2n} + C_{2n - 1}^{n} \left( \frac{1}{2} \right)^{2n - 1} + \ldots + C_{n}^{n} \left( \frac{1}{2} \right)^n = 1. 
\]
两边同乘 \(2^n\)，
\[
\left( \frac{1}{2} \right)^n C_{2n}^{n} + \left( \frac{1}{2} \right)^{n - 1} C_{2n - 1}^{n} + \ldots + C_{n}^{n} = 2^n. 
\]
即
\[
C_{n}^{0} + \frac{1}{2} C_{n + 1}^{1} + \frac{1}{4} C_{n + 2}^{2} + \frac{1}{8} C_{n + 3}^{3} + \ldots + \frac{1}{2^n} C_{2n}^{n} = 2^n. \qedhere
\]
\end{proof}


\begin{example}\label{example:例3. 6}
  利用概率论的想法证明下列恒等式:
\[
1 + \frac{A - a}{A - 1} + \frac{\left( {A - a}\right) \left( {A - a - 1}\right) }{\left( {A - 1}\right) \left( {A - 2}\right) } + \ldots  + \frac{\left( {A - a}\right) \cdots 2 \cdot  1}{\left( {A - 1}\right) \cdots \left( {a + 1}\right) a} = \frac{A}{a},
\]
其中 \(A,a\) 都是正整数,且 \(A > a\) . 
\end{example}

\begin{proof}
  设袋中有 \(A\) 个球,其中有 \(a\) 个红球,其余为白球. 从袋中随机地一个一个地取球, 摸出的球不放回. 求迟早取出红球的概率. 设事件 \(B\) 为事件 “第一次抽到红球”. 就有 B1 “表示第一次就抽到红球”, B2 表示 “第一次抽到白球,第二次抽到红球”,. . . , \(\mathrm{{Bi}}\) 表示 “前 \(i - 1\) 次都抽到了白球,第 \(i\) 次抽到了红球”. 可知最多到第 \(A - a + 1\) 次必定会取出红球. 就有 \(B = {B}_{1} + {B}_{2} + \cdots  + {B}_{A - a + 1}, {B}_{i} \cap  {B}_{j} = \varnothing ,i \neq  j\) . 

因此就有
\[
P\left( B\right)  = P\left( {\sum\limits_{{i = 1}}^{{A - a + 1}}{B}_{i}}\right)  = 1,
\]
\[
P\left( {B}_{1}\right)  = \frac{a}{A},P\left( {B}_{2}\right)  = \frac{\left( {A - a}\right) a}{A\left( {A - 1}\right) },\ldots ,P\left( {B}_{A - a + 1}\right)  = \frac{\left( {A - a}\right) \left( {A - a - 1}\right) \cdots 2 \cdot  1}{A\left( {A - 1}\right) \cdots \left( {a + 1}\right) a}. 
\]
将上述各值带入得到
\[
\frac{a}{A} + \frac{\left( {A - a}\right) a}{A\left( {A - 1}\right) } + \ldots  + \frac{\left( {A - a}\right) \left( {A - a - 1}\right) \cdots 2 \cdot  1}{A\left( {A - 1}\right) \cdots \left( {a + 1}\right) a} = 1,
\]
等式两边同乘 \(\frac{A}{a}\) 可得,
\[
1 + \frac{A - a}{A - 1} + \frac{\left( {A - a}\right) \left( {A - a - 1}\right) }{\left( {A - 1}\right) \left( {A - 2}\right) } + \ldots  + \frac{\left( {A - a}\right) \left( {A - a - 1}\right) \cdots 2 \cdot  1}{\left( {A - 1}\right) \left( {A - 2}\right) \cdots \left( {a + 1}\right) a} = \frac{A}{a}. \qedhere
\]
\end{proof}

% 在这两个例题的证明过程中, 通过巧妙地构建摸球的等可能事件概率模型, 再利用必然事件概率为 1 的性质, 从概率模型的角度简洁明了地完成了恒等式的证明, 充分展示了如何构造简单的概率模型来证明等式. 

\begin{example}
  证明 \[\sum_{i=2}^{n+1} i(i-1)C_{n+1}^i = n(n+1)2^{n-1}. \]
\end{example}

\begin{proof}
现有 \(n + 1\) 箱子排成一排，把红、黑、白三种颜色的球任意 \(n + 1\) 个放到 \(n + 1\) 个箱里，每个箱子里放一个球，且 \(n + 1\) 个红球编号为 \(1,2,\ldots,n\)，使 \(n + 1\) 个球中恰好有两个红球。

记事件 \(A_i\) 为恰好有 \(n + 1 - i\) 个白球。总事件数为 \(A_{n + 1}^{2} \cdot 2^{n - 1}\)，事件 \(A_i\) 的事件数为 \(C_{n + 1}^{n + 1 - i} \cdot A_{i}^{2}\)，于是就有
\[
P(A_i) = \frac{C_{n + 1}^{n + 1 - i} A_{i}^{2}}{A_{n + 1}^{2} 2^{n - 1}} = \frac{i \cdot (i - 1) C_{n + 1}^{i}}{n(n + 1) 2^{n - 1}}, \quad i = 2,3,\ldots,n + 1. 
\]
\(i\) 不同时，\(A_i\) 互不相容，又有
\[
\sum_{i = 2}^{n + 1} P(A_i) = 1. 
\]
即
\[
\sum_{i = 2}^{n + 1} \frac{i(i - 1) \cdot C_{n + 1}^{i}}{n(n + 1) \cdot 2^{n - 1}} = 1. 
\]
两边同乘 \(n(n + 1) 2^{n - 1}\) 后得
\[
\sum_{i = 2}^{n + 1} i \cdot (i - 1) \cdot C_{n + 1}^{i} = 1 \times 2 C_{n + 1}^{2} + 2 \times 3 C_{n + 1}^{3} + \ldots + n(n + 1) C_{n + 1}^{n + 1} = n(n + 1) 2^{n - 1}. \qedhere
\]
\end{proof}

\begin{example}
  证明 \[\sum_{i+j+k=n} \frac{n!}{i!j!k!} = 3^n,\]其中 \(i,j,k \in \{0,1,\ldots,n\}\)。
\end{example}

\begin{proof}
现设袋内有 \(3n\) 个球，其中红白黑三色球各有 \(n\) 个。记事件 \(A_{i,j,k}\) 为每次摸取一个球，摸出后放回，取 \(n\) 次，其中有 \(i\) 次摸到红球，\(j\) 次摸到白球，\(k\) 次摸到黑球。

于是有 \(P(A_{i,j,k}) = \frac{n!}{i! j! k!} \left( \frac{1}{3} \right)^i \left( \frac{1}{3} \right)^j \left( \frac{1}{3} \right)^k\)，其中 \(i,j,k \in \{0,1,2,\ldots,n\}\) 且 \(i + j + k = n\)。

又由事件 \(A_{i,j,k}\) 满足对不同取值时相互独立，得 \(\sum_{i + j + k = n} P(A_{i,j,k}) = 1\)。

整理化简
\[
\sum_{i + j + k = n} \frac{n!}{i! j! k!} \left( \frac{1}{3} \right)^{i + j + k} = \sum_{i + j + k = n} \frac{n!}{i! j! k!} \left( \frac{1}{3} \right)^n = 1. 
\]
移项后即有
\[
\sum_{i + j + k = n} \frac{n!}{i! j! k!} = 3^n, \quad i,j,k \in \{0,1,2,\ldots,n\}. 
\]
\end{proof}

\begin{example}
  证明 \(\sum_{k=0}^r C_n^k C_{n-k}^{2r-2k} 2^{2r-2k} = C_{2n}^{2r}\)。
\end{example}
\begin{proof}
现一箱子里有 \(n\) 双不同型号的鞋子，记事件 \(A_k\) 为从箱子中随机拿出 \(2r\) 只鞋子，恰有 \(k\) 双同型号的鞋子。从 \(n\) 双鞋子随机拿出 \(2r\) 只，共有 \(C_{2n}^{2r}\) 种拿法。

恰有 \(k\) 双同型号，就有 \(C_{n}^{k} C_{n - k}^{2r - 2k} 2^{2r - 2k}\) 种取法。因此就有
\[
P(A_k) = \frac{C_{n}^{k} C_{n - k}^{2r - 2k} 2^{2r - 2k}}{C_{2n}^{2r}}. 
\]
\(k\) 不同时，任意 \(A_k\) 不相容，就有
\[
\sum_{k = 0}^{r} P(A_k) = \sum_{k = 0}^{r} \frac{C_{n}^{k} C_{n - k}^{2r - 2k} 2^{2r - 2k}}{C_{2n}^{2r}} = 1. 
\]
等式两边同乘 \(C_{2n}^{2r}\) 即有
\[
\sum_{k = 0}^{r} C_{n}^{k} C_{n - k}^{2r - 2k} 2^{2r - 2k} = C_{2n}^{2r}. 
\]
\end{proof}

\section{概率模型证明恒等式方法概述}
在梳理了现有的概率模型证明中学恒等式实例后，可以发现概率模型证明中学恒等式的一般步骤是：
% （1）观察恒等式结构
% （2）代数变换，使恒等式一边为1
% （3）根据已变形的恒等式，构造合适的概率模型
% （4）结合概率论中其他重要知识梳理证明思路
% （5）证明恒等式
\begin{enumerate}
  \item 观察恒等式结构
  \item 代数变换，使恒等式一边为1
  \item 根据已变形的恒等式，构造合适的概率模型
  \item 结合概率论中其他重要知识梳理证明思路
  \item 证明恒等式
\end{enumerate}

现有的实例中，最常用的概率模型是等可能事件概率模型. 在选择实例后，要选择契合的其他重要知识来配合证明，例如实例中用到最多的必然事件概率和为1. 在现有实例中我们也可以看到，概率模型证明法不能缺少传统证明方法的辅助，概率模型法不能说是独立的新方法，在结合已有方法之后才能更加地适用. 

在实例中，本文特意将举出来两个传统法证明恒等式的例子在此证明，主要是为了与传统方法作对比展示概率模型证法的独到之处. 
在得出概率模型证明方法的基本步骤之后，再回头观察现有的例子，发现已有的例子中，概率模型证明恒等式还仅仅局限于组合恒等式和部分由组合恒等式转化的代数恒等式，而中学恒等式的三大类型还有三角不等式没有出现. 这一现实情况也说明了两个问题，一是概率模型证明法的研究没有深入，二是概率模型证明法存在很大的局限性. 

% \subsection{实例分析}
% \begin{example}
%   证明恒等式: \((a+b)^3 = {a}^{3} + 3{a}^{2}b + {3a}{b}^{2} + {b}^{3}\;\left( {a > b > 0}\right). \)
% \end{example}
% \begin{proof}
%   构建如下独立重复试验概率模型, 设进行三次独立的抽奖试验, 每次箱里面都有 \(a\) 个标记为 \(a\) 的球和 \(b\) 个标记为 \(b\) 的球(球除标记外无其他区别), 每次抽奖有放回. 设事件 \(A\) 表示在一次抽奖中抽到标记为 \(a\) 的球, 其概率 \(P\left( A\right)  = \frac{a}{a + b}\) ; 事件 \(B\) 表示在一次抽奖中抽到标记为 \(b\) 的球, 其概率 \(P\left( B\right)  = \frac{b}{a + b}\) . 

% 抽奖的结果一共有四种, 即(i)三次都抽到 a 球；(ii)恰有一次抽到 a 球；(iii)恰有两次抽到 a 球；(iiii)三次都没抽到 b 球. 

% 同 \ref{subsec:3. 1. 2} 的实例分析, 我们可以得到一个类似的等式
% \[
% \frac{{a}^{3}}{{\left( a + b\right) }^{3}} + \frac{3{a}^{2}b}{{\left( a + b\right) }^{3}} + \frac{{3a}{b}^{2}}{{\left( a + b\right) }^{3}} + \frac{{b}^{3}}{{\left( a + b\right) }^{3}} = 1. 
% \]
% 等式两边同乘 \({\left( a + b\right) }^{3}\) , 得到恒等式的证明. 
% \end{proof}


% \begin{example}
%   证明恒等式 \(\sum\limits_{{k = 0}}^{n}{\binom{n}{k} }^{2} = \binom{2n}{n}\) . 
% \end{example}

% \begin{proof}
%   构造一个如下的概率模型, 设有两组质地均匀的硬币, 分别投掷 \(n\) 次, 设事件 \(A\) 为 “两组试验中正面朝上的次数相等”. 

% 两组组试验中正面朝上 \(k\) 次的概率都为 \(\binom{n}{k} {\left( \frac{1}{2}\right) }^{n}\) . 

% 由于两组试验独立, 则正面朝上次数均为 \(k\) 的概率为
% \[
% {\left\lbrack  \binom{n}{k} {\left( \frac{1}{2}\right) }^{n}\right\rbrack  }^{2} = {\binom{n}{k} }^{2}{\left( \frac{1}{2}\right) }^{2n},
% \]
% 可得
% \[
% P\left( A\right)  = \sum\limits_{{k = 0}}^{n}{\binom{n}{k} }^{2}{\left( \frac{1}{2}\right) }^{2n} = 1. 
% \]

% 两组正面次数相等可以简化为总正面次数为任意 \(k\) , 等价于从 \(2n\) 次试验中选取 \(n\) 次正面, 组合数为 \(\binom{2n}{n}\) , 可知概率为 \(P\left( A\right)  = \binom{2n}{n} {\left( \frac{1}{2}\right) }^{2n}\) . 
% 就有
% \[
% \sum\limits_{{k = 0}}^{n}{\binom{n}{k} }^{2}{\left( \frac{1}{2}\right) }^{2n} = \binom{2n}{n} {\left( \frac{1}{2}\right) }^{2n},
% \]
% 两边同时消去 \({\left( \frac{1}{2}\right) }^{2n}\) , 得 \(\sum\limits_{{k = 0}}^{n}{\binom{n}{k} }^{2} = \binom{2n}{n}\) . 
% \end{proof}

% 投掷硬币和有放回摸球是独立重复试验概率模型最常见得两个例子. 通过这两个实例可以看出, 运用独立重复试验概率模型证明恒等式, 能够将抽象的数学式子转化为具体的试验场景, 使证明过程更加直观易懂. 



% \section{运用特殊概率模型——以巴拿赫火柴模型为例}
% \subsection{巴拿赫火柴问题}
% 特殊的概率模型往往不具有普适性, 这类概率模型适用的范围有限. 我们这里以巴拿赫火柴问题为例, 介绍这一类特殊概率模型. 巴拿赫火柴模型建立在这样一个问题上: 数学家的左右衣带中各放有一盒装有 $N$ 根火柴的火柴盒, 每次抽烟时任取一盒用一根, 求发现一盒用光时另一盒有 \(r\) 根的概率. 这个问题的概率  
% \(P = \binom{N}{2N-r} {\left( \frac{1}{2}\right) }^{{2N} - r + 1}\). 
% 由这个问题得到了一个巴拿赫恒等式 \[\sum \limits_{{r = 0}}^{n}{C}_{{2n} - r}^{n}{\left( \frac{1}{2}\right) }^{{2n} - r} = 1 . \]

% \subsection{实例分析}

% \begin{example}
%   证明恒等式 \(\mathop{\sum }\limits_{{k = 0}}^{N}\binom{N+k}{k} \frac{1}{{2}^{k}} = {2}^{N}\) . 
% \end{example}

% \begin{proof}
%   构造一个如下概率模型: 设有两个抽奖箱子, 每个箱子里都放有 \(N\) 个白球, 每次抽烟奖时取一个箱子拿出一个球, 拿出后不放回. 在发现一箱抽完时, 另一个箱子剩下 \(r\) 个球的概率为
% \[
% \binom{2N-r}{N} {2}^{-{2N} + r}. 
% \]

% 任取 \(r\) 就有
% \[
% \sum\limits_{{r = 0}}^{N}\binom{2N-r}{N} {2}^{-{2N} + r} = 1,
% \]
% 上式两边同乘 \(2N\) 并利用组合性质可得
% \[
% \sum\limits_{{r = 0}}^{N}\binom{2N-r}{N} 2^{-N + r} = {2}^{N}. 
% \]
% 令 \(k = N - r\) , 就有
% \[
% \sum\limits_{{k = N}}^{0}\binom{N+k}{k} {2}^{-k} = {2}^{N},
% \]
% 即
% \[
% \sum\limits_{{k = 0}}^{N}\binom{N+k}{k} \frac{1}{{2}^{k}} = {2}^{N}. \qedhere
% \]
% \end{proof}